\phantomsection
\section*{前言}
\addcontentsline{toc}{section}{前言}

\begin{TranslatorNoteBox}
	译者在翻译过程中进行大量了的重组及根据自己的理解进行润色。
\end{TranslatorNoteBox}

Avali 及其相关的设定和概念属于其创作者 RyuujinZERO 所有。虽然本书是基于该内容撰写的，但你不应将其视为官方设定。
本书旨在将 Ryuujin 本人已公开阐述的 Avali 原始设定与核心概念系统整理成册。此外，本书的作者、编辑、插画师及全体参与者都在已经建立的基础上仔细完善和扩展。
由于本书的特殊性质，几乎所有原始概念都保持不变，仅有两处重大修改：Avali 的核心温度从 $-30\,^\circ\mathrm{C}$ 调整为 $-7\,^\circ\mathrm{C}$，并且体液改为了氨水混合物。
这么做是为了扩大 Avali 的适宜活动温度范围，避免纯氨沸腾和纯水结冰。此外，该调整还带来了额外的优势：由于混合体液提升了整体温度，能量传递效率也因此显著提高。

% ? 此段删除了`and we worked tirelessly within theestablished parameters of the original concepts.`，因为这句话有些多余且不通顺。
本书由一群热爱 Avali 的志愿者团队无偿完成，他们承担了撰稿、插画、编辑和审校等所有环节。尽管原作者未直接参与本书撰写，但书中探讨的一切内容，均在原作者的理念与初衷下撰写。

\strikeThrough {
	本书将持续修订和完善。因此本指南讨论的主题可能会改变或删除。
	本书中的更改将以粗黑线标示（参见侧边示例），该黑线位于页面边缘，受影响段落的左侧或右侧。
	内容修订页位于目录页前。版本号从0开始按升序排列。最新版本发布后，所有先前版本都将失效。
	版本号位于每页的左上角。
	本文档有两个版本：公开版和支持者版。
	每个版本都会在版本号末尾添加一个两位字母代码，分别为“PE”或“SE”。
	草稿版本会在版本号末尾添加“D”。
	发布日期位于页面右上角。
}

如果您（读者）发现本书有问题或有适合纳入本设定指南的想法，您可以向本书的作者反馈，我们会进行处理。但这并不保证您的想法一定会被收录或认可。
未经作者、插画师或编辑同意您不允许修改任何内容。

感谢您抽出时间阅读我们的作品。我们希望您能从中获得关于本主题的有用信息和深入见解。
此外我们需要再次重申：这并非官方设定。

\clearpage
